\subsection{Applications}

PROPOSITION 2. Let A be a commutative ring, \(\mathfrak{S}\) an ideal of A and B a commutative Noetherian A-algebra such that B is contained in the Jacobson radical of B. Then every finitely generated B-module M is an ideally Hausdorff A-module with respect to \(\mathfrak{S}\) .

We shall see more generally that for every finitely generated A-module N, \(N \otimes_{A} M\) is Hausdorff with the \(\Im\) -adic topology. For \(\mathrm{N}_{(\mathrm{B})} = \mathrm{N} \otimes_{\mathrm{A}} \mathrm{B}\) is a finitely generated B-module and the B-module \(N \otimes_{A} M\) is canonically identified with \(\mathrm{N}_{(\mathrm{B})} \otimes_{\mathrm{B}} M\) by virtue of the associativity of the tensor product. Let L be the Jacobson radical of B; as \(\Im B\) is contained in L, the \(\Im\) -adic topology on \(N \otimes_{A} M\) is therefore identified with a finer topology than the \(\Im\) -adic topology on \(\mathrm{N}_{(\mathrm{B})} \otimes_{\mathrm{B}} M\) ; but this latter topology is Hausdorff since \(\mathrm{N}_{(\mathrm{B})} \otimes_{\mathrm{B}} M\) is a finitely generated B-module (no. 1, Example 1), whence the conclusion.

PROPOSITION 3. Let A be a commutative ring, B a commutative A-algebra, \(\mathfrak{I}\) an ideal of A and M a B-module. Suppose that B is a Noetherian ring and a flat A-module and that M is ideally Hausdorff with respect to \(\mathfrak{S}\) B. The following conditions are equivalent:

\begin{enumerate}
\def\labelenumi{(\alph{enumi})}
\item
  M is a flat B-module.
\item
  \(\mathbf{M}\) is a flat A-module and \(\mathbf{M} / \Im \mathbf{M} = \mathbf{M} / (\Im \mathbf{B})\mathbf{M}\) is a flat (B/3B)-module.
\end{enumerate}

If further the canonical homomorphism \(\mathbf{A} / \Im \to \mathbf{B} / \Im \mathbf{B}\) is bijective, conditions (a) and (b) are also equivalent to:

\begin{enumerate}
\def\labelenumi{(\alph{enumi})}
\setcounter{enumi}{2}
\tightlist
\item
  M is a flat A-nodule.
\end{enumerate}

Condition (a) implies (b) by Chapter I, §2, no. 7, Corollaries 2 and 3 to Proposition 8 and the fact that M/Im is isomorphic to M ⊗B (B/ImB). Suppose condition (b) holds; to show that M is a flat B-module, we shall apply Theorem 1 of no. 2 with A replaced by B and S by SIB. It will therefore be sufficient to show that the canonical mapping \(f\colon \mathfrak{B} \otimes_{\mathbf{B}} \mathbf{M} \to \mathfrak{M}\) is injective. Let \(f_1\) be the canonical mapping \(\mathfrak{I} \otimes_{\mathbf{A}} \mathbf{B} \to \mathfrak{B}\) and \(f_2\) the canonical isomorphism \(\mathfrak{I} \otimes_{\mathbf{A}} \mathbf{M} \to (\mathfrak{I} \otimes_{\mathbf{A}} \mathbf{B}) \otimes_{\mathbf{B}} \mathbf{M}\) ; \(f \circ (f_1 \circ 1_{\mathbf{M}}) \circ f_2\) is the canonical mapping \(f'\colon \mathfrak{I} \otimes_{\mathbf{A}} \mathbf{M} \to \mathfrak{IM}\) , as is easily verified. Now \(f'\) is an isomorphism since M is a flat A-module, whilst \(f_1\) is an isomorphism because B is flat over A; \(f\) is then an isomorphism.

Let \(\rho: \mathring{\mathrm{A}}/\mathfrak{I} \to \mathrm{B}/\mathfrak{I}\mathrm{B}\) be the canonical homomorphism; the (A/ \(\mathfrak{I}\) )-module structure on M/ \(\mathfrak{I}\) M derived by means of \(\rho\) from its (B/ \(\mathfrak{I}\) B)-module structure is isomorphic to that on M \(\otimes_{\mathrm{A}} (\mathrm{A}/\mathfrak{I})\) . Then it follows that, if M is a flat A-module, M/ \(\mathfrak{I}\) M is a flat (A/ \(\mathfrak{I}\) )-module and hence also a flat (B/ \(\mathfrak{I}\) B)-module if \(\rho\) is an isomorphism; we have thus proved that (c) \(\Rightarrow\) (b) in that case.

COROLLARY. Let A be a commutative Noetherian ring, \(\Im\) an ideal of A, \(\hat{\mathbf{A}}\) the Hausdorff completion of A with respect to the \(\Im\) -adic topology and M an ideally Hausdorff \(\hat{\mathbf{A}}\) -module with respect to \(\Im\hat{\mathbf{A}}\) . For M to be a flat A-module, it is necessary and sufficient that M be a flat \(\hat{\mathbf{A}}\) -module.

We know in fact that \(\hat{\mathbf{A}}\) is a Noetherian ring ( \(\S 3\) , no. 4, Proposition 8) and a flat A-module ( \(\S 3\) , no. 4, Theorem 3), that \(\Im \hat{\mathbf{A}} = \hat{\mathfrak{I}}\) ( \(\S 2\) , no. 12, Proposition 16) and that the canonical homomorphism \(\mathrm{A} / \Im \to \hat{\mathrm{A}} / \hat{\mathfrak{I}}\) is bijective ( \(\S 2\) , no. 12, Proposition 15); Proposition 3 can therefore be applied.

PROPOSITION 4. Let A and B be two commutative Noetherian rings, h: A → B a ring homomorphism, ∑ an ideal of A and ∑ an ideal of B containing ∑B and contained in the Jacobson radical of B. Let A be the Hausdorff completion of A with respect to the ∑-adic topology and B the Hausdorff completion of B with respect to the ∑-adic topology; h is continuous with these topologies and h: A → B therefore makes B into an A-algebra. Let M be a finitely generated B-module and M its Hausdorff completion with respect to the ∑-adic topology; the following properties are equivalent:

\begin{enumerate}
\def\labelenumi{(\alph{enumi})}
\item
  M is a flat A-module.
\item
  \(\hat{\mathbf{M}}\) is a flat A-module.
\item
  \(\hat{\mathbf{M}}\) is a flat \(\hat{\mathbf{A}}\) -module.
\end{enumerate}

As B with the \(\mathfrak{L}\) -adic topology is a Zariski ring, \(\hat{\mathbf{B}}\) is a faithfully flat B-module (§ 3, no. 5, Proposition 9) and \(\hat{\mathbf{M}}\) is canonically isomorphic to \(\mathbf{M} \otimes_{\mathbf{B}} \hat{\mathbf{B}}\) (§ 3, no. 4, Theorem 3); it is immediately verified that this canonical isomorphism is an isomorphism of the A-module structure on \(\hat{\mathbf{M}}\) onto the A-module structure on \(\mathbf{M} \otimes_{\mathbf{B}} \hat{\mathbf{B}}\) derived from that on M. Applying Proposition 4 of Chapter I, § 3, no. 2 with R replaced by B, S by A, E by \(\hat{\mathbf{B}}\) , F by M, we see that for M to be a flat A-module, it is necessary and sufficient that \(\hat{\mathbf{M}}\) be a flat A-module. Moreover, \(\hat{\mathbf{M}}\) is a finitely generated \(\hat{\mathbf{B}}\) -module and \(\mathfrak{J}\hat{\mathbf{B}}\) is contained in \(\hat{\mathfrak{L}} = \mathfrak{L}\hat{\mathbf{B}}\) and hence in the Jacobson radical of \(\hat{\mathbf{B}}\) (§ 3, no. 4, Proposition 8); therefore \(\hat{\mathbf{M}}\) is an ideally Hausdorff \(\hat{\mathbf{A}}\) -module with respect to \(\mathfrak{J}\hat{\mathbf{A}}\) (Proposition 2). Conditions (b) and (c) are therefore equivalent by the Corollary to Proposition 3.


